\documentclass[12pt]{article}
\usepackage[utf8]{inputenc}
\usepackage[T1]{fontenc}
\usepackage[margin=1in]{geometry}
\usepackage{setspace}
\usepackage{xcolor}
\usepackage[colorlinks=true, urlcolor=blue!60!black, linkcolor=blue!60!black]{hyperref}
\setlength{\parindent}{0em}
\setlength{\parskip}{0.6em}
\renewcommand{\baselinestretch}{1.0}
\begin{document}
\begin{center}
{\LARGE \textit{Green v. County School Board of New Kent County}, 391 U.S. 430 (1968)\par}
\vspace{0.5em}
{\large Case Brief --- Two-Page Edition\par}
\end{center}

\section{Memory Jogger}
A formerly segregated school board cannot discharge its duty to eliminate the dual system ``root and branch'' with a ``freedom-of-choice'' plan that leaves schools racially identifiable when zoning would work (\href{https://supreme.justia.com/cases/federal/us/391/430/}{Justia}).

\section{Detailed Case Facts}
Rural New Kent County, Virginia, half Black and residentially integrated, ran one white and one Black school, with buses on overlapping routes. A decade after \textit{Brown}, no one had transferred. Sued in 1965, the Board adopted freedom of choice; after three years no white student attended Watkins, and 85\% of Black students remained there.

\section{Procedural History}
The district court approved the plan, and the en banc Fourth Circuit affirmed that approval (remanding only on faculty), over the partial disagreement of Judges Sobeloff and Winter.

\section{Judicial Votes}
Vacated and remanded, 9--0. Justice Brennan wrote for the Court; there were no separate opinions.

\section{Holding}
A board that ran a state-compelled dual system must take whatever steps are necessary to convert to a unitary system, across students, faculty, staff, transportation, activities, and facilities. Its plan must work now. Freedom of choice is not unconstitutional per se but fails when it does not dismantle the dual system and better alternatives exist.

\section{Analysis of Opinions}
\textbf{Brennan:} Opening the doors ``merely begins, not ends, our inquiry.'' After eleven years of delay, ``[t]he time for mere `deliberate speed' has run out.'' Freedom of choice is only a means; in a county where geographic zoning or splitting schools by grade would integrate at once, a plan leaving 85\% of Black children at Watkins wrongly shifted the burden from the Board to Black families.

\section{Examples: Future Applications}
\textbf{Same-Side:} Integrating zones paired with free transfers that let white students return to formerly white schools perpetuate the dual system. \textbf{Opposite-Side:} Racial imbalance in a district that never practiced intentional segregation triggers no \textit{Green} duty. \textbf{Fence-Sitter:} A magnet-school choice plan raising white enrollment at formerly Black schools from 0\% to 30\% may or may not offer ``real promise.''

\section{Critique}
Conservatives say \textit{Green} turned \textit{Brown}'s colorblind rule into race-conscious ``compulsory integration'' and decades of judicial management. Progressives praise it for giving \textit{Brown} teeth but note that tying it to de jure systems enabled later retreat and resegregation.

\section{Key Quotations}
\begin{quote}
``School boards such as the respondent then operating state-compelled dual systems were nevertheless clearly charged with the affirmative duty to take whatever steps might be necessary to convert to a unitary system in which racial discrimination would be eliminated root and branch.''
\end{quote}
\noindent --- Justice Brennan, for a unanimous Court, 391 U.S. at 437--38.

\end{document}
